\documentclass[10pt,a4paper]{amsart}
\usepackage[margin=19mm]{geometry}
\usepackage{amsmath,amssymb,mathtools}
\usepackage[scr=rsfs,cal=euler]{mathalfa}
\usepackage{xfrac}
\usepackage[unicode,hidelinks,bookmarks=false]{hyperref}
\newcommand{\ie}{i.e.\ }
\newcommand{\cf}{cf.\ }
\newcommand{\resp}{respectively\ }
\newcommand{\Subsection}{\subsection}
\providecommand{\nobreakdash}{\nobreak\hbox{-}}
\setlength{\emergencystretch}{2em}
\multlinegap0pt
\usepackage{xcolor}
\usepackage{amssymb}
\usepackage{bbm}
\usepackage{esint}
\usepackage{subfig}
\usepackage{caption}
\usepackage{dsfont}
\usepackage{tikz}
\usepackage[shortlabels]{enumitem}
\usepackage[normalem]{ulem}
\newtheorem{assumption}{Assumption}
\newtheorem{definition}{Definition}
\newtheorem{theorem}{Theorem}
\newtheorem{lemma}{Lemma}
\newcommand{\N}{\mathbb{N}}
\newcommand{\PP}{\mathcal{P}}
\newcommand{\R}{{\mathbb R}}
\title{Guaranteeing Higher Order Convergence Rates for Accelerated Wasserstein gradient Flow Schemes}
\author{Raymond Chu, Matt Jacobs}
\date{Source: arXiv:2511.10884v2 (CC BY 4.0)}
\begin{document}
\maketitle

\noindent\emph{Source and licence.} Guaranteeing Higher Order Convergence Rates for Accelerated Wasserstein gradient Flow Schemes. Version arXiv:2511.10884v2 (CC BY 4.0), distributed by arXiv under the Creative Commons Attribution 4.0 International licence, http://creativecommons.org/licenses/by/4.0/, as recorded in the arXiv licence metadata for this exact version and shown on the abstract page https://arxiv.org/abs/2511.10884v2. Full licence notice: \texttt{licenses/arxiv-2511.10884-CC-BY-4.0.txt}.\par\smallskip

\noindent\textit{Editorial scope.} Counted unit: the article's lemma estimates\_lambda\_negative (source0.tex lines 2247--2264). The article's complete original proof of it is delivered with it (source0.tex lines 2266--2336, one \texttt{proof} environment of 71 lines). No mathematics is written, repaired or re-derived: the statement and the proof are the article's own lines, in the article's own notation, and no retained line is substituted or re-flowed.  Retained uncounted as the source-local context the proof needs: lines 324--330; lines 335--345; lines 347--354; lines 348--348; lines 362--393; lines 557--576; lines 1471--1476; lines 1997--2027; lines 2091--2097; lines 2102--2110; lines 2112--2125; lines 2146--2151.  Nothing was omitted from any retained span, so the package carries no omission record.  No retained line contains a citation, so the wrapper carries no bibliography and no bibliography entry.  No retained line contains a figure environment, an image-inclusion command or drawing code, so no image is delivered and the package declares no asset dependency.  Editorial formatting only, in addition: an AMS \texttt{amsart} preamble that restates the article's own declarations and macros used by the retained lines, namely the environments \texttt{assumption}, \texttt{definition}, \texttt{theorem}, \texttt{lemma} on separate counters, together with the standard packages those lines need.  The retained environments print in this document as Assumption 1, Definition 1, Theorem 1, Lemma 1 in order of appearance, whereas the article prints the same items with the numbering of its own full document.  The complete native source of the article as distributed in the arXiv e-print for this exact version is bundled unchanged as \texttt{sources/189074/source0.tex}.
\begin{equation}
\label{gradient_flow_eqn12}
\begin{cases}
 \partial_t \rho(t,x) - \nabla\cdot \Bigl(\rho(t,x) \nabla\Bigl(\tfrac{\delta \phi}{\delta \mu}\bigl(\rho(t,\cdot),x\bigr)\Bigr)\Bigr) = 0 \quad \text{on } (0,\infty)\times\mathbb{R}^d,\\[6pt]
 \rho(0,\cdot) = \rho_0(\cdot) \quad \text{on } \mathbb{R}^d,
\end{cases}
\end{equation} can be interpreted formally as the gradient flow of $\phi$ on the Wasserstein space $(\PP_2(\R^d), \mathbf{W}_2)$.

To construct our scheme, we exploit both the Eulerian (Wasserstein) and Lagrangian $L^2$ gradient flow structures of the equation to construct a novel second-order-accurate-in-time numerical scheme for solving \eqref{gradient_flow_eqn12}. To motivate this Lagrangian $L^2$ perspective, let us assume that the velocity field $v=-\nabla \bigl( \frac{\delta \phi}{\delta \mu} \bigr)$ is sufficiently regular, so that the solution $\rho(t,x)$ can be expressed as the push-forward of the initial data $\rho_0\in \PP_2(\R^d)$ by the Lagrangian flow $X(t,x)$, i.e.
\begin{equation}
\rho_t(x) = (X(t,\cdot)_{\#} \rho_0)(x). \label{pushforward_rho_def}
\end{equation}
where
\begin{equation}
\begin{cases} 
\displaystyle \frac{d}{dt}X(t,x)  =-\nabla \left[ \frac{\delta \phi}{\delta \mu}( X(t,\cdot)_{\#} \rho_0,x) \right] \quad \text{on } [0,\infty),  \\
X(0,x) = x, \label{characteristic_eqn}
\end{cases}
\end{equation}

As we will see below, the Lagrangian flow \eqref{characteristic_eqn} is a gradient flow of the lifted energy functional $X\mapsto \phi^{\#}_{\rho_0}(X)$ over the Hilbert space $L^2(\R^d;\rho_0)$ where
\begin{equation} \phi^{\#}_{\rho_0}(X) := \phi(X_{\#} \rho_0). \label{lift_intro_def} \end{equation} That is, if $\nabla \phi^{\#}_{\rho_0}(X)$ denotes the $L^2(\R^d;\rho_0)$ Fréchet derivative of the map $X \mapsto \phi^{\#}_{\rho_0}(X)$, then
\begin{equation}
\begin{cases} 
\displaystyle \frac{d}{dt}X(t,x) = -\nabla \phi^{\#}_{\rho_0}(X(t,x)) \quad \text{on } [0,\infty),  \\
X(0,\cdot) = \textnormal{Id}. \label{L2_Lagrangian_Flow}
\end{cases}.
\end{equation}

\begin{equation} \phi^{\#}_{\rho_0}(X) := \phi(X_{\#} \rho_0).  \end{equation} That is, if $\nabla \phi^{\#}_{\rho_0}(X)$ denotes the $L^2(\R^d;\rho_0)$ Fréchet derivative of the map $X \mapsto \phi^{\#}_{\rho_0}(X)$, then

To obtain a second-order-in-time discretization of \eqref{gradient_flow_eqn12}, we apply the second order trapezoidal finite difference scheme to \eqref{L2_Lagrangian_Flow}. In particular, we consider
\begin{equation}
\begin{cases}
X^{\tau}_{n+1}
\in \operatorname*{arg\,min}_{\xi \in L^2(\mathbb{R}^d;\rho_0)}
\left[
\frac{1}{2}\bigl(\phi^{\#}_{\rho_0}(\xi)
+ \langle
\nabla \phi^{\#}_{\rho_0}(X^{\tau}_n),
\xi
\rangle_{L^2(\mathbb{R}^d;\rho_0)}
\bigr)
+ \frac{1}{2 \tau}\|\xi - X^{\tau}_n\|_{L^2(\mathbb{R}^d;\rho_0)}^2
\right], \\
\rho^{\tau}_{n+1} := (X^{\tau}_{n+1})_{\#} \rho_0,
\end{cases}
\label{trapezoid_variational_form}
\end{equation}
where the initial condition $X^{\tau}_0$ is prescribed.
 The optimizers $X^{\tau}_{n+1}$ of \eqref{trapezoid_variational_form} admits the following trapezoidal finite difference scheme:
\begin{equation}
X^{\tau}_{n+1}
= X^{\tau}_n
- \frac{\tau}{2} \bigl(
\nabla \phi^{\#}_{\rho_0}(X^{\tau}_{n+1})
+ \nabla \phi^{\#}_{\rho_0}(X^{\tau}_{n})
\bigr). \label{trapezoid_rule123}
\end{equation}
We note that, unlike the JKO scheme, the trapezoidal update maps
$T^{\tau}_{n+1} := X^{\tau}_{n+1} \circ (X^{\tau}_n)^{-1}$
are not, in general, optimal transport maps between
$\rho_{n}^{\tau}$ and $\rho_{n+1}^{\tau}$, since they need not be gradients of convex functions.

\begin{assumption} \label{assume_rough}
    We fix a measure $\rho_0 \in \mathcal{P}_2(\mathbb{R}^d)$ and let $\mathbb{H} := L^2(\mathbb{R}^d;\rho_0)$.  \smallskip 

\noindent Given a functional $\phi : \mathcal{P}_2(\mathbb{R}^d) \to \mathbb{R}$, we assume that its lifted functional $\phi_{\rho_0}^{\#} : \mathbb{H} \to \mathbb{R}$, defined by \eqref{lift_intro_def}
is Fréchet differentiable over $\mathbb{H}$ and is $\lambda$-convex over $\mathbb{H}$ for some $\lambda \in \mathbb{R}$ Additionally, we assume that the lifted energy functional is proper, i.e.,
\[
\inf_{\xi \in L^2(\mathbb{R}^d;\rho_0)} \phi^{\#}_{\rho_0}(\xi) > -\infty.
\] 

\noindent We further assume that in \eqref{trapezoid_variational_form} the time step parameter satisfies
\[
    \lambda/2 + 1/\tau > 0.
\]
Finally, we suppose that the initial iterates satisfy
\[
    \lim_{\tau \downarrow 0} \|X_0^{\tau} - \mathrm{Id}\|_{\mathbb{H}} = 0,
    \qquad
    \sup_{\tau > 0} \|\nabla \phi^{\#}_{\rho_0}(X_0^{\tau})\|_{\mathbb{H}} < \infty.
\]
\end{assumption}

\begin{equation}
    W_2(\rho_n^{\tau},\rho_m^{\tau}) 
    \leq \|X^{\tau}_{n}-X^{\tau}_m\|_{\mathbb{H}} 
    \leq    \tau \cdot |n-m| \cdot \sqrt{\tilde{\mathcal{C}}(\lambda,\tau,T)}  \cdot \|\nabla \phi^{\#}_{\rho_0}(X^{\tau}_0)\|_{\mathbb{H}}.
    \label{lipschitz_bound}
\end{equation}

\begin{definition}[Interpolations] \label{interpolations_usual}
Fix a time step size $\tau > 0$. Then we define
\[
\ell_{\tau}(t) := \frac{t - n\tau}{\tau}, \quad
\overline{X}^{\tau}_t := X^{\tau}_{n+1},  \text{ and }
\underline{X}^{\tau}_t := X^{\tau}_n, \quad \text{for } t \in [n\tau,(n+1)\tau).
\]
The linearly interpolated numerical Lagrangian flow and energy functional are given by
\[
\tilde{X}^{\tau}_t := (1 - \ell_{\tau}(t)) \underline{X}^{\tau}_t + \ell_{\tau}(t) \overline{X}^{\tau}_t, \quad
\phi_{\tau}(t) := (1 - \ell_{\tau}(t)) \phi^{\#}_{\rho_0}(\underline{X}^{\tau}_t) 
+ \ell_{\tau}(t) \phi^{\#}_{\rho_0}(\overline{X}^{\tau}_t).
\]
For $\xi_1,\xi_2 \in \mathbb{H}$, we define the metric
\[
d(\xi_1, \xi_2) := \|\xi_1 - \xi_2\|_{\mathbb{H}},
\]
and for any $\xi \in \mathbb{H}$, the linearly interpolated distance functional is
\[
d^2_{\tau}(t;\xi) := 
(1 - \ell_{\tau}(t)) d^2(\xi, \underline{X}^{\tau}_t) 
+ \ell_{\tau}(t) d^2(\xi, \overline{X}^{\tau}_t).
\] 
Define the constant interpolation gradient error term as
\[
G^{\tau}_t := \|\nabla \phi^{\#}_{\rho_0}(X^{\tau}_n)\|^2_{\mathbb{H}} - \|\nabla \phi^{\#}_{\rho_0}(X^{\tau}_{n+1})\|^2_{\mathbb{H}}, \quad t \in [n\tau,(n+1)\tau).
\] and the distance between iterates as
\[
\mathcal{D}_{t}^{\tau} := d(\underline{X}^{\tau}_t,\overline{X}^{\tau}_t).
\]
\end{definition}

\begin{theorem}[Approximate EVI] \label{differential_inequality_thm}
Under Assumption \ref{assume_rough} and the notation introduced therein, as well as the notation from Definition \ref{interpolations_usual}, the following differential inequalities holds: 
\begin{equation}
   \frac{d}{dt} d^2_{\tau}(t;\xi) + \lambda d^2_{\tau}(t;\xi) 
    \leq  2(\phi^{\#}_{\rho_0}(\xi) - \phi_{\tau}(t)) + \frac{\tau}{2} G^{\tau}_t + 2\mathcal{R}_t^{\tau} + |\lambda|  (\mathcal{D}_t^{\tau}) \cdot((\mathcal{D}_t^{\tau}) + d_{\tau}(t;\xi)).
\end{equation}
\end{theorem}

\begin{definition} \label{further_interpolated_d}
Using the notation from Definition \ref{interpolations_usual}, let $\tau, \eta > 0$ be two time steps. We define the further interpolated distance functional as
\small{\[
d^2_{\tau,\eta}(t;s) := (1-\ell_{\tau}(t)) [  (1-\ell_{\eta}(s)) \cdot d^2(\underline{X}^{\eta}_s,\underline{X}^{\tau}_t) + \ell_{\eta}(s) \cdot d^2(\overline{X}^{\eta}_s,\underline{X}^{\tau}_t) ]
 +  \ell_{\tau}(t) \left( (1-\ell_{\eta}(s)) d^2(\underline{X}^{\eta}_s,\overline{X}^{\tau}_t) + \ell_{\eta}(s) d^2(\overline{X}^{\eta}_s,\overline{X}^{\tau}_t) \right) \]}
\[ = (1-\ell_{\eta}(s)) \cdot d^2_{\tau}(t; \underline{X}^{\eta}_s) + \ell_{\eta}(s) \cdot d^2_{\tau}(t;\overline{X}^{\eta}_s). \]


\end{definition}

\begin{theorem}[Differential Inequality]  \label{differential_inequality_different_time_steps}
Fix two time steps $\tau, \eta > 0$ satisfying
\[
\frac{\lambda}{2} + \min\left\{\frac{1}{\tau},\frac{1}{\eta}\right\} > 0.
\]
Using the assumptions and notation from Theorem \ref{differential_inequality_thm} and from Definition \ref{further_interpolated_d}, we have
\[
\frac{d}{dt} d^2_{\tau,\eta}(t;t) + 2 \lambda d^2_{\tau,\eta}(t;t) 
\leq \frac{\tau}{2} G^{\tau}_t + \frac{\eta}{2} G^{\eta}_t
+ 2(\mathcal{R}^{\tau}_t + \mathcal{R}^{\eta}_t) 
+ |\lambda| \left( (\mathcal{D}_{t}^{\tau})^2 + (\mathcal{D}_{t}^{\eta})^2 \right) 
+ |\lambda| (\mathcal{D}_t^{\tau} + \mathcal{D}_t^{\eta}) d_{\tau,\eta}(t;t).
\]
\end{theorem}

\begin{lemma}[Bounds on $\mathcal{R}^{\tau}_t$] \label{R_bounds}
Under Assumption \ref{assume_rough} and the notation of Theorem \ref{differential_inequality_thm}, the following bound holds:
\[
\mathcal{R}_t^{\tau} \leq \left( 1-\ell_{\tau}(t) \right) \left( \frac{\tau}{4} G^{\tau}_t - \frac{\lambda}{2} (\mathcal{D}_t^{\tau})^2 \right) - \frac{1}{\tau}\left( \ell_{\tau}(t) - \frac{1}{2} \right) (\mathcal{D}_t^{\tau})^2.
\]
\end{lemma}

\begin{lemma} \label{estimates_lambda_negative}
Under the notation and assumptions of Theorem \ref{differential_inequality_different_time_steps}, and with the additional assumptions that $\lambda < 0$, we have the following estimates for any $T \geq 0$:
\begin{equation}
\begin{aligned}
    \int_0^T e^{2 \lambda_{\tau} t} (\mathcal{D}_t^{\tau})^2 \,dt 
    &\leq \tau^2 (T+\tau) e^{-2 \lambda_{\tau} \tau} \|\nabla \phi^{\#}_{\rho_0}(X_0^{\tau})\|^2_{\mathbb{H}},  
    \label{distance_squared_bound} \\
    \int_0^T e^{2 \lambda_{\tau} t} G_t^{\tau} \,dt 
    &\leq \tau \|\nabla \phi^{\#}_{\rho_0}(X_0^{\tau})\|^2_{\mathbb{H}}, \\
    \int_0^T e^{2 \lambda_{\tau} t} \mathcal{R}^{\tau}_t \,dt 
    &\leq \tau^2\left( \frac{1}{8} + e^{-2 \lambda_{\tau} \tau} \left( \frac{1}{2} - \frac{2}{3} \lambda_{\tau}(T+\tau) \right) \right) 
    \|\nabla \phi^{\#}_{\rho_0}(X^{\tau}_0)\|^2_{\mathbb{H}}, \\ 
    \int_0^T e^{\lambda_{\tau} t} (\mathcal{D}_t^{\tau}) \,dt 
    &\leq \tau (T+\tau) e^{-\lambda_{\tau} \tau} \|\nabla \phi^{\#}_{\rho_0}(X_0^{\tau})\|_{\mathbb{H}}.
\end{aligned}
\end{equation}

\end{lemma}

\begin{proof}

Choose the smallest integer $N$ such that $T \leq N \tau$. \smallskip 

\noindent \textbf{Step 1: Bounding } $\displaystyle \int_0^T e^{2 \lambda_{\tau} t} (\mathcal{D}_t^{\tau})^2\,dt$.

We first observe that
\begin{equation}
\int_0^T e^{2\lambda_{\tau} t} d^2(\underline{X}_t^{\tau},\overline{X}^{\tau}_t)\,dt 
\leq \int_0^{N \tau} e^{2\lambda_{\tau} t} d^2(\underline{X}_t^{\tau},\overline{X}^{\tau}_t)\,dt 
= \underbrace{\frac{e^{2\lambda_{\tau} \tau} - 1}{2 \lambda_{\tau}}}_{\geq 0} \sum_{j=0}^{N-1} e^{2 \lambda_{\tau} j \tau} d^2(X^{\tau}_j,X^{\tau}_{j+1}).
\label{first_integral_boundss}
\end{equation} 

\noindent By applying \eqref{lipschitz_bound}, we obtain
\begin{equation}
d^2(X^{\tau}_j,X^{\tau}_{j+1}) \leq e^{-2\lambda_{\tau} (j+1) \tau} \tau^2 \|\nabla \phi^{\#}_{\rho_0}(X_0^{\tau})\|^2_{\mathbb{H}}. \label{distance_bounds123}
\end{equation}
Substituting this into \eqref{first_integral_boundss} and using the inequalities $(e^{2 \lambda_{\tau} \tau} - 1) \geq 2\lambda_{\tau} \tau$ and $\lambda_{\tau} \leq 0$, we obtain the desired bound. \smallskip

\noindent \textbf{Step 2: Bounding $\int_0^T e^{2 \lambda_{\tau}t} G^{\tau}_t dt$}  \smallskip

Because $G^{\tau}_t$ is piecewise constant on $[n \tau,(n+1)\tau)$, we obtain:
\[
\int_0^T e^{2 \lambda_{\tau}t}  G^{\tau}_s ds \leq \max \left\{ \int_0^{N \tau}  e^{2 \lambda_{\tau}t}  G^{\tau}_s ds , \int_0^{(N-1)\tau}  e^{2 \lambda_{\tau}t}  G^{\tau}_s ds  \right\} 
\]  
\[
= \max \left\{ \sum_{j=0}^{N} G^{\tau}_{j \tau} \int_{j \tau}^{(j+1)\tau} e^{2 \lambda_{\tau}t}  ds, \sum_{j=0}^{N-1} G^{\tau}_{j \tau} \int_{j \tau}^{(j+1)\tau} e^{2 \lambda_{\tau}t}  ds  \right\} =  \frac{e^{2 \lambda_{\tau} \tau} - 1}{2 \lambda_{\tau}}  \left( \sum_{j=0}^{N-1} G_{j \tau}^{\tau}  e^{2 \lambda_{\tau} j \tau} + e^{2 \lambda_{\tau} N\tau}(G^{\tau}_{N\tau})_{+} \right).  
\] Where $(x)_+ := \max\{x,0\}$. To simplify our notations for this proof, we define for $n \in \N$
\begin{equation}
\alpha_n := ||\nabla \phi^{\#}_{\rho_0}(X_n^{\tau})||^2_{\mathbb{H}}, \label{alpha_def}
\end{equation}
so that from Definition \ref{interpolations_usual} $G^{\tau}_{j \tau} = \alpha_j - \alpha_{j+1}$. Since $\lambda_{\tau} < 0$, we obtain for any $m \in \N$
\[
\sum_{j=0}^{m} G^{\tau}_{j \tau} e^{2 j \lambda_{\tau} \tau} = \sum_{j=0}^{m} (\alpha_j - \alpha_{j+1}) e^{2 j \lambda_{\tau} \tau} = \alpha_0 + \sum_{j=1}^{m} \underbrace{(e^{2 j \lambda_{\tau} \tau} - e^{2 (j-1) \lambda_{\tau} \tau})}_{\leq 0} \alpha_j - e^{2m \lambda_{\tau} \tau} \alpha_{m+1} \leq \alpha_0.
\]  
Thus, we conclude from $(e^{2\lambda_{\tau} \tau} - 1)/(2 \lambda_{\tau}) >0$ that we have the bound
\[ \frac{e^{2 \lambda_{\tau} \tau} - 1}{2 \lambda_{\tau}}  \left( \sum_{j=0}^{N-1} G_{j \tau}^{\tau}  e^{2 \lambda_{\tau} j \tau} + e^{2 \lambda_{\tau} N\tau}(G^{\tau}_{N\tau})_{+} \right) \leq \frac{e^{2 \lambda_{\tau} \tau} - 1}{2 \lambda_{\tau}} ||\nabla \phi^{\#}_{\rho_0}(X_0^{\tau})||^2_{\mathbb{H}} \leq \tau ||\nabla \phi^{\#}_{\rho_0}(X_0^{\tau})||^2_{\mathbb{H}} .
\]

\noindent \textbf{Step 3: Bounding $\int_0^t e^{2 \lambda_{\tau} t} \mathcal{R}_t^{\tau} dt$}

By Lemma \ref{R_bounds}, we have that
\[ \int_0^t e^{2 \lambda_{\tau} t} \mathcal{R}^{\tau}_t dt \leq \frac{\tau}{4}\int_0^t e^{2 \lambda_{\tau} t}  (1-\ell_{\tau}(t)) \cdot G^{\tau}_t dt + \int_0^t e^{2 \lambda_{\tau} t}\left( -\frac{\lambda_{\tau}}{2}  (1-\ell_{\tau}(t)) - \frac{1}{\tau}(\ell_{\tau}(t)-\frac{1}{2}) \right)  (\mathcal{D}_{t}^{\tau})^2) dt.  \]
\[ = (I) + (II). \] We first focus on bounding $(II)$. First observe that because $\lambda_{\tau} < 0$ and $j \in \N$,
\[ \int_{j \tau}^{(j+1)\tau}  \frac{e^{2 \lambda_{\tau} t}}{\tau}(\frac{1}{2}-\ell_{\tau}(t)) dt = - \frac{ (\lambda_{\tau} \tau + e^{2 \lambda_{\tau} \tau}(\lambda_{\tau} \tau - 1) + 1)e^{2 \lambda_{\tau} j \tau} }{4\lambda_{\tau}^2 \tau^2} \leq -\frac{\lambda_{\tau} \tau}{6} e^{2 \lambda_{\tau} j \tau}.  \]  In the above inequality, we used $-(x+e^{2x}(x-1)+1) \leq -2x^3/3$. Hence, we have from \eqref{distance_bounds123}
\[ \int_0^{(N-1)\tau}  \frac{e^{2 \lambda_{\tau} t}}{\tau} (\frac{1}{2} - \ell_{\tau}(t)) \cdot d^2(\overline{X}^{\tau}_t,\underline{X}_t^{\tau}) dt \leq -\frac{\lambda_{\tau} \tau}{6} \sum_{j=0}^{N-1} e^{2 \lambda_{\tau} j \tau} \cdot d^2(X^{\tau}_j,X^{\tau}_{j+1}) \leq -\frac{\lambda_{\tau} \tau^2 }{6}   e^{-2 \lambda_{\tau} \tau} \cdot T \cdot ||\nabla \phi^{\#}_{\rho_0}(X^{\tau}_0)||^2_{\mathbb{H}}. \] We also have from \eqref{lipschitz_bound}
\[ \int_{(N-1) \tau}^t \frac{e^{2 \lambda_{\tau} t}}{\tau} (\frac{1}{2} - \ell_{\tau}(t)) \cdot (\mathcal{D}_t^{\tau})^2 dt \leq \frac{1}{2}  e^{2 \lambda_{\tau}(N-1) \tau}d^2(X^{\tau}_{N-1},X^{\tau}_{N}) \leq  \frac{\tau^2}{2} e^{-2 \lambda_{\tau} \tau} \cdot ||\nabla \phi^{\#}_{\rho_0}(X_0^{\tau})||^2_{\mathbb{H}} .     \]


Now for the first integral term in $(II)$, we have from Step 1
\[ \int_0^t e^{2\lambda_{\tau} t} \left( -\frac{\lambda_{\tau}}{2}(1-\ell_{\tau}(t)) \right) d^2(\overline{X}_{t}^{\tau},\underline{X}_t^{\tau}) \leq -\frac{\lambda_{\tau}}{2} \int_0^t e^{2\lambda_{\tau} t} d^2(\underline{X}_t^{\tau},\overline{X}^{\tau}_t) dt \leq -\frac{\lambda_{\tau}}{2} \tau^2(T+\tau)e^{-2 \lambda_{\tau} \tau} ||\nabla \phi^{\#}_{\rho_0}(X_0^{\tau})||^2_{\mathbb{H}}, \] where the final inequality is due to our Step 1 bound. \smallskip

Combining these bounds gives us
\[ (II) \leq  \tau^2e^{-2\lambda_{\tau} \tau} \left( \frac{1}{2} -  \frac{2}{3} \lambda_{\tau}(T+\tau)   \right) \cdot||\nabla \phi^{\#}_{\rho_0}(X^{\tau}_0)||^2_{\mathbb{H}}. \] 

Now we focus on $(I)$. First recall the notation $\alpha_n$ from \eqref{alpha_def}. By using that $G^{\tau}_{t}$ is piece-wise constant on $[n\tau,(n+1)\tau)$ and $(1-\ell_{\tau}(t)) \geq 0$ the arguments of Step 1 imply
\begin{equation}
(I) \leq \frac{\tau}{4} \left( \sum_{j=0}^{N-1} (\alpha_n-\alpha_{n-1}) \cdot \int_{j \tau}^{(j+1)\tau}  e^{2 \lambda_{\tau} s} (1-\ell_{\tau}(s)) ds + (G^{\tau}_{N \tau})_{+} \cdot \int_{(N-1)\tau}^{N \tau} e^{2 \lambda_{\tau} s} (1-\ell_{\tau}(s)) ds    \right). \label{bound_on_I_1}
\end{equation} By integrating the term $e^{2 \lambda_{\tau} s} (1-\ell_{\tau}(s))$, we see that for any $m \in \N$
\[ \sum_{j=1}^{m} (\alpha_n-\alpha_{n-1}) \cdot \int_{j \tau}^{(j+1)\tau}  e^{2 \lambda_{\tau} s} (1-\ell_{\tau}(s)) ds = 
\underbrace{\frac{ e^{2 \lambda_{\tau} \tau} - (1+2\lambda_{\tau} \tau)  }{4 \lambda_{\tau}^2 \tau}}_{\geq 0} \sum_{j=1}^{m} (\alpha_j-\alpha_{j-1}) e^{2 \lambda_{\tau} j \tau}.  \] By a similar argument as in summation by parts step of Step 2, we have that
\[ \frac{ e^{2 \lambda_{\tau} \tau} - (1+2\lambda_{\tau} \tau)  }{4 \lambda_{\tau}^2 \tau} \sum_{j=1}^{m} (\alpha_j-\alpha_{j-1}) e^{2 \lambda_{\tau} j \tau} \leq \frac{ e^{2 \lambda_{\tau} \tau} - (1+2\lambda_{\tau} \tau)  }{4 \lambda_{\tau}^2 \tau}  \cdot ||\nabla \phi^{\#}_{\rho_0}(X_0^{\tau})||^2_{\mathbb{H}}.  \]
Because $\lambda_{\tau} \tau < 0$, we have that $e^{2\lambda_{\tau} \tau} - (1+2\lambda_{\tau}  \tau) \leq 2\lambda_{\tau}^2 \tau^2$, so we have for any $m \in \N$
\[
\sum_{j=0}^{m} (\alpha_j-\alpha_{j+1}) \int_{j \tau}^{(j+1)\tau} e^{2 \lambda s} (1-\ell_{\tau}(s)) ds \leq \frac{\tau}{2} \cdot ||\nabla_{\mathbf{W}} \phi(\rho^{\tau}_0,X^{\tau}_0)||^2_{\mathbb{H}}.
\]  Therefore, we obtain from \eqref{bound_on_I_1} that
\[ (I) \leq \frac{\tau^2}{8} \cdot ||\nabla_{\mathbf{W}} \phi(\rho^{\tau}_0,X^{\tau}_0)||^2_{\mathbb{H}} .  \] Therefore, by combining our bounds on $(I)$ and $(II)$, we conclude the third step. \smallskip

\noindent \textbf{Step 4: Bounding $\int_0^t e^{\lambda_{\tau} t} (\mathcal{D}_t^{\tau}) dt$}
The desired bound follows from Step 1 and Holder's inequality.\end{proof}

\end{document}
